\section{Imperfections}

\begin{frame}{Recap: Assumptions}
    \mycomment{
    So far, we have a nice edifice of theory building on pretty rigid assumptions. Let us see what changes if we drop them, one by one.
    }
    \mytext{
    \begin{itemize}
        \item no market entry barriers
        \item perfect information
    \end{itemize}
    }
\end{frame}

\begin{frame}{Recap: Assumptions}
    \mytext{
    \begin{itemize}
        \item perfect capital markets:
        \item[] if idea is profitable, funding will follow
        \item no individual discrimination
        \begin{itemize}
            \item by race, age, gender
        \end{itemize}
        \item right to sell and buy
        \begin{itemize}
            \item could be restricted
            \item feudal obligations
            \item other qualifications
        \end{itemize}
        \item right to property
        \item[] $\rightarrow$ for ever?
    \end{itemize}
    }
\end{frame}

\begin{frame}{Patents, Intellectual Property}
    \mytext{
    A patent can be granted on
    \begin{itemize}
        \item a solution
        \item to a concrete technical problem
    \end{itemize}
    You cannot get a patent on
    \begin{itemize}
        \item ideas
        \item mathematical insights
        \item rules
    \end{itemize}
    }
    \mycomment{
    The legal protection for intellectual property is different between the EU (written law) and the Common Law countries.

    A chief difference is that between (engl.) \textit{copyright}---the right to produce copies, that you can sell---and the unsellable\textit{ droit d'auteur} (Urheberrecht). 

    Generally, existing business is protected; a patent not used is difficult to defend. 

    Registration of trade marks and commercial patterns protects from fast movers to use an idea (or pattern) before the inventor does.
    }
\end{frame}

\begin{frame}{Otherwise Protected Assets}
    \mycomment{
    There are plenty other reasons why market entry could be difficult or impossible. Since businesses are typically eager to protect their market, many sort of (often illegal) barriers to entry can be devised.
    }
    \mytext{
        \begin{itemize}
            \item Factor Availability
            \item Macroeconomic Circumstance
            \item Military secrets
            \item Big Data
            \item Algorithms
            \item Capacity restrictions
            \item Access to buyers
            \item Contracts with intermediate product providers
        \end{itemize}
    }
\end{frame}

\begin{frame}{Imperfect Competition: Quasi-Monopoly}
    \mycomment{
    Markets are not perfect. The owner of a little kebab shop competes with the shop down the road, but only to an extent. On his road, she has a quasi-monopoly.
    
    As long as she is not massively more expensive than the market price, her neighbours will still go to her. But at some point, they will.

    If, on the contrary, she sets her price cheaper than the neighbour, at some point people will travel from far.

    So the demand curve she is actually facing has the shape of an inverted S. There is an inelastic part in the middle where she can set the price rather freely, while there are highly elastic parts on top and bottom.
    }
    \showmytikz[1]{../shared/tikz/tikz-imperfect-comp-pricesetting.tex}
\end{frame}

\begin{frame}{Imperfect Competition: Quasi-Monopoly}
    \mycomment{
    You see that the largest profit is realised at point $C$. You would set the price at that point so high just before everyone starts walking away. 

    The extra surplus is a quasi-monopoly surplus. Notice that our graph shows a relatively large inelastic part of the demand curve. In reality, that part will likely be short. 
    
    Also note that the aggregated demand curve $DD$ usually sits somewhere much further out, and that our quasi-monopolists supply curve (her average costs) may not be exactly identical to aggregated supply $SS$.
    }
    \showmytikz[2]{../shared/tikz/tikz-imperfect-comp-pricesetting.tex}
\end{frame}

\begin{frame}{Asymmetric Information}
    \mycomment{
    One assumption was that everyone always knows everything about any product. But this is often not the case.  
    
    Caveat emptor---buyer, beware! The buyer normally knows less about the product than the seller. This leads to market imperfections.
    }
    \mytext{
    Information is a key issue on labour markets, insurance markets, and capital markets:
    \begin{itemize}
        \item Moral Hazard
        \item[] insured people are less careful
        \item Adverse Selection
        \item[] only not good things need advertising
        \item[] only risk seekers seek insurance
        \item Principal-Agent-Problem
        \item[] Agent knows more than the principal
        \item[] incentives differ
        \item[] think client, architect, builder
    \end{itemize}
    }
\end{frame}

\begin{frame}{Market for Lemons}
    \mycomment{
    Information asymmetries can have long-term effects on markets, which distorts prices and also distorts what is on offer. 
    
    A classic example from \noun{Akerlof} is that of used cars. Imagine I took perfectly care of the car I now want to sell, with gears and motor treated well. I have no real option to proof that to my buyers. Others offer the same model at a lower price. Am I likely to sell my car?
    }
    \mytext{
    `Lemon' is slang for a bad car sold as good.
    \begin{itemize}
        \item buyers less informed than sellers
        \item buyers tend to go for cheapest offer
        \item bad offers push out good offers
        \item long-term quality hard to know
        \item competition via price and features, not quality
        \item planned obsolescence, fast-moving consumer goods
        \item sellers of quality are pushed out
    \end{itemize}
    }    
\end{frame}

\begin{frame}{Risk and Uncertainty}
    \mycomment{
    In reality, many trades are not instantaneous, and are prone to risk. Will the wares arrive in time? Will the price be the same once they arrive? 

    We can differentiate between risk and uncertainty (Frank \noun{Knight});
    \begin{itemize}
        \item risk: insurable
        \item uncertainty: non-insurable
    \end{itemize}
    In practice, we can sort the risks identified by impact and likelihood and come up with strategies to mitigate, hedge, insure etc. against them. Perfect insurance markets should make risk disappear.
    }
    \showmytikz{../shared/tikz/tikz-riskregister}
\end{frame}

\begin{frame}{Risk and Uncertainty}
    \mycomment{
    This scheme is known as Rumsfeld-Matrix (after a US Secretary of Defence who popularised it). There are things that we know, and things we don't know. But are we actually aware of what we know and what we don't know?
    }
    \showmytikz{../shared/tikz/tikz-rumsfeld}
\end{frame}

\begin{frame}{Tragedy of the Commons}
    \begin{tabular}{l|l|l|}
                    & rival & non-rival  \\\hline
         excludable & private goods & club goods \\\hline
         non-excludable & common goods & public good \\\hline
    \end{tabular}
\end{frame}

\begin{frame}{Tragedy of the Commons}
    \mytext{
    \begin{itemize}
        \item Over-use of public good
        \item Externality of over-use
        \item Game theory
        \item privatisation, or administration
    \end{itemize}
    }
\end{frame}


\begin{frame}{Tragedy of the Commons}
    \showmytikz[1]{../shared/tikz/tikz-tragedyofthecommons.tex}
    \mycomment{
        "Commons" are (traditionally) lands everyone is allowed to use, such as a pond or pool in the middle of a village, surrounding forests used for firewood, etc.
    
        The problem is that while profits I get from the commons are private (the fish I fish), the costs are communal (to feed the fish and clean the pool). 

        Since I can profit without paying, and since I don't profit if I pay (or take care), there is a strong incentive to over-use the commons, and take too little care. 

        Imagine an apple tree in the middle of the village: it is likely that the apples will be harvested too early.

    }
\end{frame}

\begin{frame}{Coase Theorem}
    \mytext{
    If there are no transaction costs, then
    \begin{itemize}
        \item bargaining always leads to efficient outcomes
        \item efficient outcome independent of initial distribution
        \item sound property rights must be guaranteed
        \item citizens must be able to sue for social damage
    \end{itemize}
    }
    \mycomment{
    The Coase Theorem is often misunderstood, or cited too excessively. Transaction costs are never zero. Property rights can be well defined, but might be difficult to defend. While it is possible to sue, this comes at sometimes high costs. And the initial distribution need to still allow for some bargaining power.

    The argument was significantly improved by Elinor \noun{Ostrom}, showing that communities can reach consensus on the use of common pools without privatisation and without central management, with much less strict assumptions than \noun{Coase}.
    }
\end{frame}


\begin{frame}{External Effects}
    \showmytikz[1]{../shared/tikz/tikz-externalities-pigou}
    \mycomment{
    Whatever we do, we typically incur external effects. We are not paying for all consequences of our actions; we are using common goods without paying; and even if we are forced to pay, some damages are not redeemable by payment.
    }
\end{frame}

\begin{frame}{External Effects}
    \showmytikz[1,2]{../shared/tikz/tikz-externalities-pigou}
    \mycomment{
    Quite often, the costs of the producer do not cover all costs. Those costs not covered are called externalities, or social costs.

    A classic example is the pollution caused by production of the good, or in fact any other type of damage or costs.

    These are hard to grasp, and are often only realised much in the past or by future generations. We abstract from those problems and assume $MC_{social}$ is known.
    }
\end{frame}

\begin{frame}{External Effects}
    \showmytikz[1,2,3]{../shared/tikz/tikz-externalities-pigou}
    \mycomment{
        Social costs are above the costs covered by the producers and the price the consumer pay them. Some of them are offset by the consumer and producer surplus: yes, we destroy the planet, but hey, we consume!

        Bear in mind that the consumer realising the surplus may not be the same consumers suffering the damage!

        However, there is an excess damage that is neither translated into consumer nor producer surplus. This is in the literature called dead weight loss, but it is actually excess costs, or non-covered damage. (Dead weight loss is welfare we could have but do not have due to market imperfections. Here, we actually have a loss, so it is not "dead weight").
    }
\end{frame}

\begin{frame}{External Effects}
    \showmytikz[1,2,3,4]{../shared/tikz/tikz-externalities-pigou}
    \mycomment{
        The idea could be to internalise the social costs, lifting the cost curve (supply curve) up to meet the social cost curve. Then the price realised would be realistic. This could be done by a tax.

        This tax can achieve two things at once:
        \begin{itemize}
            \item shifting the price so that private costs equal social costs
            \item raising a government revenue that can be used to compensate those who suffer without benefiting.
        \end{itemize}
    }
\end{frame}

